\subsection{Tensor Products of Simple Modules}

Let \(A_{1}\) and \(A_{2}\) be algebras over the commutative field K. We denote the K-algebra \(A_{1} \otimes A_{2}\) by A.

Lemma 1. --- Let \(M_{1}\) and \(N_{1}\) be \(A_{1}\) -modules, and let \(M_{2}\) and \(N_{2}\) be \(A_{2}\) -modules. We make the following assumptions:

\begin{enumerate}
\def\labelenumi{(\roman{enumi})}
\item
  The \(A_{1}\) -module \(M_{1}\) is finitely generated.
\item
  The \(A_{2}\) -module \(M_{2}\) is finitely generated, or \(N_{1}\) is finite-dimensional over K.
\end{enumerate}

Set \(\mathrm{M} = \mathrm{M}_1\otimes \mathrm{M}_2\) and \(\mathrm{N} = \mathrm{N}_1\otimes \mathrm{N}_2\) , and view them as modules over the ring \(\mathrm{A} = \mathrm{A}_1\otimes \mathrm{A}_2\) . The canonical homomorphism (II, §3, No.~5, p.~251)

\[
\lambda : \operatorname{Hom} _ {\mathrm{K}} (\mathrm{M} _ {1}, \mathrm{N} _ {1}) \otimes \operatorname{Hom} _ {\mathrm{K}} (\mathrm{M} _ {2}, \mathrm{N} _ {2}) \longrightarrow \operatorname{Hom} _ {\mathrm{K}} (\mathrm{M}, \mathrm{N})
\]

then induces an isomorphism of K-vector spaces

\[
\varphi : \operatorname{Hom} _ {\mathrm{A} _ {1}} (\mathrm{M} _ {1}, \mathrm{N} _ {1}) \otimes \operatorname{Hom} _ {\mathrm{A} _ {2}} (\mathrm{M} _ {2}, \mathrm{N} _ {2}) \longrightarrow \operatorname{Hom} _ {\mathrm{A}} (\mathrm{M}, \mathrm{N}).
\]

The mapping \(\lambda\) is injective (II, §7, No.~7, p.~308, Proposition 16) and sends the linear subspace \(\mathrm{Hom}_{\mathrm{A}_1}(\mathrm{M}_1,\mathrm{N}_1)\otimes \mathrm{Hom}_{\mathrm{A}_2}(\mathrm{M}_2,\mathrm{N}_2)\) to \(\mathrm{Hom}_{\mathrm{A}}(\mathrm{M},\mathrm{N})\) . It therefore suffices to prove that every A-linear mapping from M to N belongs to the image of \(\mathrm{Hom}_{\mathrm{A}_1}(\mathrm{M}_1,\mathrm{N}_1)\otimes \mathrm{Hom}_{\mathrm{A}_2}(\mathrm{M}_2,\mathrm{N}_2)\) by \(\lambda\) . Let \(u:\mathrm{M}\to \mathrm{N}\) be an A-linear mapping. Let \(x\in \mathrm{M}_1\) . Denote by \(u_{x}\) the \(\mathrm{A}_2\) -linear mapping \(y\mapsto u(x\otimes y)\) from \(\mathrm{M}_2\) to \(\mathrm{N}_1\otimes \mathrm{N}_2\) . Set \(\mathrm{P} = \mathrm{Hom}_{\mathrm{A}_2}(\mathrm{M}_2,\mathrm{N}_2)\) , and denote by \(\nu\) the canonical homomorphism from \(\mathrm{N}_1\otimes \mathrm{P}\) to \(\mathrm{Hom}_{\mathrm{A}_2}(\mathrm{M}_2,\mathrm{N}_1\otimes \mathrm{N}_2)\) (II, §4, No.~2, p.~269). This mapping is injective (II, §4, No.~2, p.~269, Proposition 2, (i) applied to the K-vector space \(\mathbf{N}_1\) ). By assumption (ii), there exists a linear subspace \(\mathrm{V}_x\) of \(\mathbf{N}_1\) , finite-dimensional over K, such that \(u_{x}\) takes on values in \(\mathrm{V}_x\otimes \mathrm{N}_2\) . It follows that \(u_{x}\) is the image by \(\nu\) of a unique element \(v_{x}\) of \(\mathbf{N}_1\otimes \mathbf{P}\) . The mapping \(\widetilde{u}:x\mapsto v_x\) from \(\mathbf{M}_1\) to \(\mathbf{N}_1\otimes \mathbf{P}\) is \(\mathbf{A}_1\) -linear. By assumption (i), the \(\mathbf{A}_1\) -module \(\mathbf{M}_1\) is finitely generated. A reasoning analogous to the above shows that \(\widetilde{u}\) belongs to the image of \(\mathrm{Hom}_{\mathrm{A}_1}(\mathrm{M}_1,\mathrm{N}_1)\otimes \mathrm{P}\) in \(\mathrm{Hom}_{\mathrm{A}_1}(\mathrm{M}_1,\mathrm{N}_1\otimes \mathrm{P})\) . Lemma 1 follows.

THEOREM 2. --- Let \(A_{1}\) and \(A_{2}\) be algebras over the commutative field K; let \(S_{1}\) be a simple \(A_{1}\) -module and \(S_{2}\) a simple \(A_{2}\) -module. Let \(D_{1}\) and \(D_{2}\) be the respective commutants of \(S_{1}\) and \(S_{2}\) . Set \(M = S_{1} \otimes S_{2}\) , \(A = A_{1} \otimes A_{2}\) , and \(D = D_{1} \otimes D_{2}\) . We view M as a left (A,D)-bimodule.

\begin{enumerate}
\def\labelenumi{\alph{enumi})}
\item
  The commutant of the A-module M is equal to \(\mathrm{D}_{\mathrm{M}}\) .
\item
  The mapping \(a \mapsto aM\) is an isomorphism from the set of right ideals of D, ordered by inclusion, to the set of A-submodules of M, ordered by inclusion. The inverse mapping sends a submodule N of M to the ideal of D consisting of the elements d such that \(dM \subset N\) .
\end{enumerate}

Assertion a) follows from Lemma 1 because a simple module is monogenous.

Let \(\mathrm{T}\) be the \((\mathrm{A}_1, \mathrm{D}_2)\) -bimodule \(\mathrm{S}_1 \otimes (\mathrm{D}_2)_d\) . We identify \(\mathrm{M} = \mathrm{S}_1 \otimes \mathrm{S}_2\) with \(\mathrm{T} \otimes_{\mathrm{D}_2} \mathrm{S}_2\) (II, §3, No.~8, p.~258); this identification is compatible with the structures of left modules over the ring \(\mathrm{A} = \mathrm{A}_1 \otimes \mathrm{A}_2\) .

Let \(\mathrm{N}\) be an A-submodule of \(\mathrm{M}\) ; it is an \(\mathrm{A}_2\) -submodule of \(\mathrm{T} \otimes_{\mathrm{D}_2} \mathrm{S}_2\) , stable by the endomorphisms of the form \((a_1)_\mathrm{T} \otimes 1_{\mathrm{S}_2}\) for \(a_1\) running through \(\mathrm{A}_1\) . It follows from Corollary 2 of VIII, p.~63 that there exists a unique \((\mathrm{A}_1, \mathrm{D}_2)\) -sub-bimodule \(\mathrm{V}\) of \(\mathrm{T}\) such that \(\mathrm{N} = \mathrm{V} \otimes_{\mathrm{D}_2} \mathrm{S}_2\) .

The isomorphism \(u\) from \(\mathrm{T} = \mathrm{S}_1 \otimes (\mathrm{D}_2)_d\) to \(((\mathrm{D}_1)_d \otimes (\mathrm{D}_2)_d) \otimes_{\mathrm{D}_1} \mathrm{S}_1\) defined by \(u(s \otimes d) = 1 \otimes d \otimes s\) is \((\mathrm{A}_1, \mathrm{D}_2)\) -linear. We identify these \((\mathrm{A}_1, \mathrm{D}_2)\) -bimodules. A reasoning analogous to that given above proves the existence and uniqueness of a right \((\mathrm{D}_1 \otimes \mathrm{D}_2)\) -submodule \(\mathfrak{a}\) of \(\mathrm{D}_1 \otimes \mathrm{D}_2\) such that \(\mathrm{V} = \mathfrak{a} \otimes_{\mathrm{D}_1} \mathrm{S}_1\) . In view of the identifications made above, \(\mathfrak{a}\) is the unique right ideal of \(\mathrm{D} = \mathrm{D}_1 \otimes \mathrm{D}_2\) such that \(\mathrm{N} = \mathfrak{a}\mathrm{M}\) .

We have just proved that the mapping \(a \mapsto aM\) is bijective; the last assertion follows.

COROLLARY 1. --- The module \(S_{1} \otimes S_{2}\) over the ring \(A_{1} \otimes A_{2}\) is semisimple (resp. isotypical, simple) if and only if the ring \(D = D_{1} \otimes D_{2}\) is semisimple (resp. simple, a field). In particular, \(S_{1} \otimes S_{2}\) is simple if the commutant of \(S_{1}\) or \(S_{2}\) is equal to K.

By Theorem 2, the module \(S_{1} \otimes S_{2}\) over the ring \(D_{1} \otimes D_{2}\) is semisimple (resp. isotypical, simple) if and only if the right D-module \((\mathrm{D}_{1} \otimes \mathrm{D}_{2})_{d}\) is (VIII, p.~109, Proposition 10). Now the right D-module \(D_{d}\) is simple if and only if D is a field; it is isotypical (resp. semisimple) if and only if the ring D is simple (resp. semisimple) (VIII, p.~120, Definition 1; VIII, p.~121, Corollary 1; and VIII, p.~137, Proposition 2).

COROLLARY 2. --- We have \(\mathfrak{R}_{\mathrm{A}}(\mathrm{M}) = \mathfrak{R}(\mathrm{D})\mathrm{M}\) . The A-module M is without radical if and only if the ring D is without radical.

This follows from Proposition 8 of VIII, p.~108 and Theorem 2, b).

